\chapter{Halo finding and merger histories}\label{ch:halo-finding}
\chapterstatus{written}{2026-11-14}

Halos are found with friends-of-friends (standard and relaxed), their spherical-overdensity properties measured around
the FoF centres, substructure resolved with a phase-space finder, and halos followed between snapshots by particle
overlap. All of them run in \file{SimulationCode} on the fly and in the snapshot analyser (\cref{ch:snapshot-analyzer})
through the analysis layer.

\begin{modulebox}{FriendsOfFriends, PhaseSpaceHaloFinder, HaloHistory}
\nsp{FML::FOF}: \file{FriendsOfFriends/FoF.h} (\code{FriendsOfFriends}), \file{FriendsOfFriends/FoFBinning.h}
(\cls{FoFHalo}), \file{FriendsOfFriends/FoFDomain.h}, \file{FriendsOfFriends/RelaxedFoF.h}
(\code{RelaxedFriendsOfFriends}), \file{FriendsOfFriends/SphericalOverdensity.h} (\code{SphericalOverdensity}).
\nsp{FML::PHASE_SPACE_HALOS}: \file{PhaseSpaceHaloFinder/PhaseSpaceHaloFinder.h} (\code{find}), \file{Distributed.h}
(\code{find_distributed}), \file{BindingPotential.h}, \file{Adapters.h}.
\nsp{FML::HALO_HISTORY}: \file{HaloHistory/HaloHistory.h} (\code{match}, \code{build}), \file{HaloHistory/CatalogueIO.h}.\par
Documentation: \file{FriendsOfFriends/RelaxedFoF.md}, \file{FriendsOfFriends/SphericalOverdensity.md},
\file{PhaseSpaceHaloFinder/README.md} and \file{tests/RESULTS.md}, \file{HaloHistory/README.md}.
Tests: \file{tests/test_fof.cpp}, \file{tests/test_fof_domain.cpp}, \file{tests/test_so.cpp}, \file{tests/test_so_multislab.cpp},
\file{FriendsOfFriends/tests}, \file{PhaseSpaceHaloFinder/tests}, \file{HaloHistory/tests}.
\end{modulebox}

\section{Friends-of-friends}\label{sec:halo-fof}

Two particles are friends if they are closer than $b\,\bar n^{-1/3}$, $b$ the linking length in units of the mean particle
separation (typically 0.2), and a group is a connected set of friends \cite{Davis1985}. The particles are binned to a grid of
cells larger than two linking lengths (at most about one particle per cell; the result does not depend on the grid), so only
$2^D$ neighbouring cells need to be searched, and groups are grown by recursively tagging the friends of friends. The particles
are sorted by one coordinate first (20--30\% faster).

\paragraph*{MPI.} Each task receives copies of the particles within a buffer of its slab (or pencil, \file{FoFDomain.h}, with up to 8
neighbours) from the neighbouring tasks, links its particles and the copies, and keeps the groups whose centre it owns; the FoF IDs of the
copies are sent back to their owners, and IDs are unique across tasks. The buffer must exceed the largest halo
(\param{fof_buffer_length_mpch}, default $3\Mpch$); groups that reach its edge are counted and reported, as they might be cut and found
twice. Tasks without particles take part normally.

\paragraph*{Halo properties.} The halo class is a template parameter: \cls{FoFHalo} accumulates the number of particles, the mass (sum
of \code{get_mass}), the centre of mass with periodic unwrapping, the mean velocity and the velocity dispersion; a derived class can bin
anything else. FoF IDs are stored in the particles if they have \code{set_fofid}.

\begin{cppcode}[FoF halos]
#include <FML/FriendsOfFriends/FoF.h>
std::vector<FML::FOF::FoFHalo<Particle, 3>> halos;
FML::FOF::FriendsOfFriends<Particle, 3>(part.get_particles_ptr(), part.get_npart(),
                                        0.2,      // linking length b
                                        20,       // minimum particles per halo
                                        true,     // periodic
                                        3.0 / Boxsize, halos);
\end{cppcode}

In \file{SimulationCode}: \param{fof}, \param{fof_linking_length}, \param{fof_nmin_per_halo}, \param{fof_buffer_length_mpch}, the
output format and the mass function binning. With separate CDM and baryon particles, linking lengths and particle counts refer to one
species.

\section{Relaxed FoF}\label{sec:halo-rfof}

In low-resolution (PM/COLA) simulations small halos are too diffuse for $b=0.2$ and fragment, so the mass function is too low and the
bias of small halos too high. The relaxed FoF of Dai et al.\ \cite{Dai2019} (as in FastPM \cite{Feng2016}) uses a longer linking length
for small halos and rejects unbound clumps joined by it through their velocity dispersion:
\begin{enumerate}
\item halos are binned by particle number with edges $N_i=20,40,80,160,320,\infty$;
\item the linking length of bin $i$ interpolates linearly between $\ell_1(z)=l_1-a_1/(1+z)$ for the smallest and
      $\ell_n(z)=\max[l_6-a_2/(1+z),\ell_{\min}]$ for the largest halos;
\item a halo is fake if $\sigma_v\ge r_0(z)V_{\rm std}(M,z)$, with the 3D dispersion $\sigma_v$, $r_0=b_1-b_2\ln(1+z)$ and
      $V_{\rm std}=V_0[E(z)M/10^{15}M_\odot h^{-1}]^{1/3}$, $V_0=1100$\,km/s \cite{Evrard2008};
\item pass $i$ runs FoF with $\ell_i$ on the remaining particles, saves the genuine groups smaller than $N_i$ (with at least
      \code{nmin} particles) and removes their particles and all ungrouped ones; the rest go to the next pass.
\end{enumerate}
The \code{fastpm} mode (default) stops after the $n$ passes; the \code{paper} mode continues with linking lengths shrinking by 0.9 per
pass and saves every group that passes the velocity cut (Algorithm 1 of the paper). The defaults are FastPM's
($l_1=0.25$, $a_1=0.012$, $l_6=0.24$, $a_2=0.06$, $\ell_{\min}=0.2$, $b_1=7.02$, $b_2=6.025$, $n_{\min}=8$), fitted to TNG300-2-Dark for
FastPM with a $0.3\Mpch$ mean separation and 40 steps; for other set-ups they should be refitted. $r_0(z)\le0$ for $z>2.2$, where every halo
would be rejected: the code warns and \code{r0_min} sets a floor. In \file{SimulationCode}: \param{fof_relaxed} and the
\code{fof_relaxed_*} parameters; the catalogue then has the pass, bin, linking length and $\sigma_v/V_{\rm std}$ of each halo.

\begin{validation}[Validation: FoF]
\file{tests/test_fof.cpp}: groups, centres and stored IDs equal a serial FoF for any number of tasks; buffer-edge groups are reported; custom
halo classes; the linking length in 2D. \file{tests/test_relaxed_fof.cpp} on 1, 2 and 4 tasks: members, pass, bin and size of every halo
equal an independent SciPy KD-tree implementation on a catalogue with compact, diffuse and fake clumps and cold cores in hot envelopes, in both
modes and at several redshifts; one bin without the velocity cut is standard FoF; tasks without particles.
\end{validation}

\section{Spherical overdensity}\label{sec:halo-so}

\code{SphericalOverdensity} measures halos around given centres (e.g.\ the FoF centres) using all particles, not only the group members:
\begin{enumerate}
\item the centre is refined by a shrinking sphere \cite{Power2003}: move to the centre of mass of the particles inside and shrink the radius by
      10\% until fewer than $\max(n_{\min},10\%)$ of the starting particles remain (at most 200 iterations);
\item for each threshold $\Delta$ (in units of the mean density) the radius is the \emph{outermost} radius where the mean enclosed density is
      still above $\Delta\bar\rho$: with the particles sorted by distance, the largest $i$ with $M_i/(\frac43\pi r_i^3)\ge\Delta\bar\rho$, and
      $R_\Delta$ from $M_i/(\frac43\pi R_\Delta^3)=\Delta\bar\rho$. Growing from the centre and stopping at the first crossing is not well
      defined at simulation resolution: the innermost particles are Poisson noise around a centre that is not on a particle;
\item mass-weighted mean velocity and dispersion inside the first radius.
\end{enumerate}
The definitions in \file{SimulationCode} (\param{fof_so_definitions}) are $\Delta$ times the mean (\code{200m}), the critical density
(\code{500c}) or the virial overdensity of Bryan and Norman (\code{vir}) \cite{BryanNorman1998}. The search radius is the same on all tasks
so results do not depend on the decomposition; particles come from every slab that intersects the search spheres, the buffer is refreshed
collectively if a centre moves, and halos whose sphere did not reach the threshold are flagged as truncated.

\begin{validation}[Validation: spherical overdensity]
Singular isothermal spheres on a uniform background (also on task boundaries) give the analytic radius, mass and centre for any number of
tasks; spheres beyond the search radius are flagged; cored halos of $\sim100$ particles (as in PM simulations) do not collapse to a few
particles. \file{tests/test_so_multislab.cpp}: exact ghost membership from non-adjacent slabs, variable-size records, bounded messages, empty
and non-uniform slabs, and a direct full-catalogue reference for centres, radii, masses, velocities and truncation.
\end{validation}

\section{Phase-space halos}\label{sec:halo-phasespace}

\file{PhaseSpaceHaloFinder} is an independent implementation of a hierarchical phase-space halo finder inspired by ROCKSTAR
\cite{Behroozi2013}:
\begin{itemize}
\item spatial FoF parents with linking parameter 0.28;
\item in each parent, positions and velocities normalised by the mass-weighted dispersions, a 6D FoF radius from the 70\% quantile of the
      nearest-neighbour distances, and recursive refinement into children down to a minimum core size;
\item the leaf groups seed halos; seeds consistent within the Poisson errors of their core positions and velocities are merged, and a sampled
      6D density-saddle test removes weaker peaks without a resolved dip towards a denser one;
\item particles are assigned on the way back up with a dynamical position scale from the seed's circular velocity and a reference density
      (200 times the mean by default);
\item unbinding with the full isolated Newtonian potential of the assigned halo (a monopole tree with opening angle 0.7, the spline-softened
      Green function of \cref{ch:short-range}, an optional Hubble-flow term), iterated until stable.
\end{itemize}
The output has exclusive bound memberships by particle ID, bound and inclusive masses, core position and velocity, dispersions,
$V_{\max}$ and $R_{\max}$, the bound extent (not an SO radius) and the host. It runs serially, with OpenMP and on MPI slabs: FoF chains across
any number of slabs are stitched by propagating minimum IDs to a fixed point, and complete parents are routed to a hash owner, without
gathering the catalogue. Adapters give seeds for the SO measurement, HOD halos and history snapshots. This is a first implementation, not a
calibrated ROCKSTAR replacement.

\begin{validation}[Validation: phase-space finder]
Exact potentials against direct sums; two overlapping populations separated in velocity are recovered; invariance under permutation, rotation,
boosts, translation and unit rescaling; Plummer spheres give one object with $V_{\max}$ within 1--4\% of the analytic value; identical serial
and distributed memberships. Planted subhalos on circular orbits in a Plummer host were recovered in 40 of 48 realisations (misses within
1.5--2.5 scale radii of the host centre). An isolated 4000-particle Plummer sphere gives $V_{\max}=97.10$\,km/s against ROCKSTAR's 97.05\,km/s and
a bound mass within 0.3\% of ROCKSTAR's $M_{200b}$ (\file{tests/RESULTS.md}).
\end{validation}

\section{Merger histories}\label{sec:halo-history}

\file{HaloHistory.h} links halo catalogues of successive snapshots by shared particle IDs \cite{Srisawat2013,Elahi2019}. For halos $a$ and $b$
with $N_a$, $N_b$ members of which $N_{ab}$ are shared, the merit is
\begin{equation}
\mathcal M_{ab}=\frac{N_{ab}^2}{N_aN_b}=\frac{N_{ab}}{N_a}\frac{N_{ab}}{N_b},
\end{equation}
and each halo's descendant is the candidate with the largest merit (then the most shared particles, then the smallest ID), above optional
thresholds; all candidates are kept for inspection. Sorted particle-ID indices are merge-joined, so only pairs that share a particle get a
counter. A target's main progenitor is chosen among the sources whose descendant it is. With a gap allowance a halo without a descendant is
matched to the earliest later snapshot where it reappears. Catalogues carry required metadata (mass definition, ID namespace, units) that must
agree, exclusive or inclusive membership, and an acyclic host forest; \file{CatalogueIO.h} reads and writes a versioned ASCII format and
converts to HOD halos. Matching two snapshots of $10^7$ halo particles in $10^5$ halos takes about 7\,s on one core. It is a baseline, not a
reproduction of TreeFrog, SubLink or Consistent Trees.
